% r03_fill_K5c: proof of (K_0), part VII (continued): the hybrid lower bound, the polynomial lower bound and the floor shift.

\begin{proposition}[the hybrid lower bound]\label{rot:prop:H4}
For the hybrid flow on $T$ and every real smooth periodic source $s$,
\[
\Zun_{\Xk,T}(s)\ge e^{-C(1+k)^p\delta_r^\sigma|T|}\int\chi_r\,e^{-\hs^{(r)}_T}.
\]
\end{proposition}
\begin{proof}
This is the analogue of \cite[Prop.~4.1, (4.5)]{K2} for the hybrid flow (the print has ``$+C$'' in the exponent; in a lower bound it must be ``$-C$''; the factor $(1+k)^p$ is that of \Cref{rot:cor:OF} for the hybrid flow and is harmless, $(1+k)^p\delta_r^\sigma\to0$ by \Cref{rot:lem:r}). Its inputs are: the block-spin algebra of the tile average (\Cref{rot:lem:tile}, \Cref{rot:lem:pair}); the hybrid analogues of \cite[Props.~2.1--2.3]{Ba1}, the decay and regularity of $G_j$, $\Delta^{(j)}$, $C^{(j)}$ along the hybrid flow with the paired contours, with \cite[Prop.~5.6]{K2} for the restricted inverses on unions of big blocks, which is \Cref{rot:prop:BD}(a) at every level ($j\ge2$ by its proof, $j\in\{0,1\}$ by the same proofs at the tile and first $Y$-frame levels); the cumulant lower bound for the step Gaussian, \Cref{rot:thm:R1}, whose extra lower-bound term $4n|\Lambda|e^{-\gamma_0\gamma_K^8b^2/(2n)}$ at the threshold $b_j=p_1(t_j)=b_1(1+\ln t_j^{-1})^{p_1}$, $t_j=5^j\eps$, $p_1\ge1$ \cite[(3.43), (3.50)]{Ba1}, on $\Lambda=T^{(j)}_1$ with $|T^{(j)}_1|=|T|t_j^{-3}$, is $\le C_\kappa t_j^\kappa|T^{(j)}_1|$ for every $\kappa$, of the form of Ba{\l}aban's own step error, and summed over $j<r$ is $\le C'_\sigma\delta_r^\sigma|T|$, uniformly in $T$, $k$, $r$; and Ba{\l}aban's lower-bound argument \cite[\S\S3--4]{Ba1} at block factor $5$ (\Cref{thm:a06:BALA}) on the torus on which the steps are performed. The hybrid flow runs step $1$ in $X$'s frame and steps $2,\dots,r$ in $Y$'s frame. On $T_1$, aligned for $X$ and skew for $Y$, steps $2,\dots,r$ are steps on a torus skew in the frame of their blocks, and the lower-bound argument is the hybrid instance of \Cref{rot:prop:SK}, with the constant $C(1+k)^p$ of \Cref{rot:cor:OF} for the hybrid flow (Ba{\l}aban's Neumann cube step being replaced by the constant-background reduction of \Cref{rot:lem:CB}). On $T_2$, skew for $X$ and aligned for $Y$, only step $1$ is skew: the fine axis lattice $\Xk$ on the torus with period lattice $R^{-1}P_J$; its blocks tile $T_2$ (\Cref{rot:lem:adm}(ii), \Cref{rot:lem:tile}(i)), its fluctuation covariance has mass of order $\eps^{-1}$ in physical units, so the torus images of its kernels are $O(e^{-c\,50MJ/\eps})$, and the step lands on $Y_{k-1}\subset T_2$, aligned for every later step. Both instances pass through the bridge \Cref{rot:prop:CT} (King adopts the free-boundary counterterm on \cite[p.~330]{K2} and appeals to \cite{Ba1} for (4.5)) and are stated for the centred-block hybrid action (\Cref{rot:conv:C}).
\end{proof}

\begin{proposition}[the polynomial lower bound for the rotated lattice]\label{rot:prop:H5}
For $\Yk$ on $T$ and every source $s=(zG,h_0)$ with $h_0=-n\psi$, $n\ge0$: $\Zun_{\Yk,T}(s)\ge e^{-C_1|T|}\eps^{\,q(s)}$ for some $q(s)<\infty$.
\end{proposition}
\begin{proof}
At $h=0$ this is the lower half of \cite[(4.2)]{K2} (``The results of [Ba 1--4] apply to both these models''), Ba{\l}aban's lower bound \cite[(1.14)]{Ba1}: on $T_2$, aligned in $Y$'s frame, by \Cref{thm:a06:T1,thm:a06:BALA} after the passive rotation, with the bridge \Cref{rot:prop:CT}; on $T_1$ by the skew instance \Cref{rot:prop:SK}. The source $zG$ is removed by Jensen's inequality: $\Zun_{\Yk,T}(zG,h)=\Zun_{\Yk,T}(0,h)\langle e^{zF(G)}\rangle_{(T,h),\Yk}\ge\Zun_{\Yk,T}(0,h)\,e^{z\langle F(G)\rangle}=\Zun_{\Yk,T}(0,h)$ for real $z$, since $\langle F(G)\rangle=0$ under the charge conjugation of \Cref{rot:def:tori}. At the mass-increasing wall $h_0=-n\psi$, $n>0$, the two-sided $\eps$-uniform bound \cite[(4.2)]{K2} is false (\Cref{rot:rem:log}) and only the polynomial form is asked for; it is \Cref{rot:thm:FS}, at every King point, or alternatively \Cref{rot:prop:J} from the source-free instance and an upper bound on the one-point Wick function.
\end{proof}

\begin{proposition}[Jensen: the wall instance from the source-free instance]\label{rot:prop:J}
Let $h_0=-n\psi\le0$ be a wall. For every real $t$, $\Lambda\in\{\Xk,\Yk\}$ and $T\in\{T_1,T_2\}$,
\[
\Zun_{\Lambda,T}(tG,h_0)\ \ge\ \Zun_{\Lambda,T}(0,0)\,\exp\Big[-n\sum_x\eps^3\psi(x)\big\langle|\phi(x)|^2-2\sW_\eps\big\rangle_{T,0,\Lambda}\Big],
\]
the expectation in the source-free periodic state. On $\Yk$ the one-point function $w_k(T)\defeq\langle|\phi(x)|^2-2\sW_\eps\rangle_{T,0,\Yk}$ is independent of $x$ ($P_J\subset\Yk$, \Cref{rot:lem:adm}), so the right side is $\Zun_{\Yk,T}(0,0)e^{-n\sS_{\Yk}(\psi)w_k(T)}$, and an upper bound $w_k(T)\le C(1+k)$ gives the polynomial form of \Cref{rot:prop:H5} at the wall with $q(tG,h_0)=q(0,0)+nC\sS(\psi)/\ln5$.
\end{proposition}
\begin{proof}
$W_\Lambda(h_0)=\sum_x\eps^3h_0(x)|\phi(x)|^2-2\sW_\eps\sS_\Lambda(h_0)$ gives $\Zun(tG,h_0)=e^{-2\sW_\eps\sS(h_0)}\Zun(0,0)\langle e^{tF(G)+\sum\eps^3h_0|\phi|^2}\rangle_{T,0,\Lambda}$. Put $X\defeq-tF(G)+n\sum_x\eps^3\psi(x)|\phi(x)|^2$; then $\langle X\rangle=n\sum_x\eps^3\psi(x)\langle|\phi(x)|^2\rangle$ because $\langle F(G)\rangle=0$ by charge conjugation. Jensen's inequality $\langle e^{-X}\rangle\ge e^{-\langle X\rangle}$ and $-2\sW_\eps\sS(h_0)-\langle X\rangle=-n\sum_x\eps^3\psi(x)(\langle|\phi(x)|^2\rangle-2\sW_\eps)$ give the bound. Only an upper bound on the one-point function is used, only on $\Yk$ and only without sources; the mass-decreasing sign would need a lower bound, which Jensen's inequality does not give.
\end{proof}

\begin{theorem}[floor shift: the polynomial lower bound at the wall at every King point]\label{rot:thm:FS}
Let $m^2$ be any King point and $h$ a real $P_J$-periodic function (it may depend on $k$) with $\|h\|_\infty\le H$ uniformly in $k$. Write $a_s\defeq s+\dm(s;\eps)$ for the bare scalar coefficient at the King point $s$ and let $M^2$ be a King point with $\Delta_0(M^2)\defeq\liminf_{\eps\to0}[a_M-a_m]>2H$; such an $M^2$ exists for every coupling and every $H$, since $\Delta_0(M^2)\to+\infty$ as $M^2\to\infty$ (the coefficients of $\partial_s\dm$ are non-negative at order $\lambda$, $O(s^{-1/2})$ at order $e^2$ and $O(s^{-1})$ at the two-loop orders, \cref{app:a05:B}). Then on $\Lambda\in\{\Yk,\Xk\}$, $T\in\{T_1,T_2\}$,
\[
\ln\Zun^{m^2}_{\Lambda,T}(0,h)\ge-C|T|-C(1+k)|T|-qk\ln5,\qquad\text{hence}\qquad\Zun^{m^2}_{\Lambda,T}(tG,h)\ge e^{-C_1|T|}\eps^{\,q'}\quad\text{for every real }t .
\]
This is the polynomial form of \Cref{rot:prop:H5} at the wall $h_0=-n\psi$ for every $n>0$; the sign of $h$ is used nowhere.
\end{theorem}
\begin{proof}
The inputs are: (i) the source-free lower bound \cite[(1.14)]{Ba1} at the one fixed King point $M^2$ (\Cref{thm:a06:T1}); (ii) the upper bound \cite[(1.3)]{Ba2} at one fixed $M_+^2$ with $\liminf_{\eps\to0}[a_{M_+}-a_M]>0$ (which exists by the same growth); on $T_1$ the skew instances of \Cref{rot:prop:SK}; (iii) the bounds on the coefficients of the one-point Wick function through order $13$ in the couplings on a compact mass interval, from Reisz's power-counting theorem \cite[Thm.~1]{Rei88} with the counterterm conditions \cite[(3.33)]{K1} (the order-$e^2$ coefficient is computed in \cref{app:a05:A}), and the boundedness of $Y\defeq1+\partial_s\dm$ on $[m^2,M_+^2]$ (\cref{app:a05:B}); (iv) exact identities: the bare-mass form $\Zun^{s}(T;0,u\mathbf 1)=e^{-P_s-2c_su|T|}\hat Z(a_s-2u)$ with $\hat Z(\gamma)$ the convention-free integral at bare coefficient $\gamma$ and $P_s=E_0(s)+E_1(s)$, the integrated counterterm identity for $P_M-P_m$, Jensen's inequality, the monotonicity of $\langle\sum_x\eps^3|\phi|^2\rangle_\gamma$ in $\gamma$, translation invariance on $\Lambda$ and charge conjugation.
\emph{Step 1 (floor shift).} $e^{\sum\eps^3h|\phi|^2}=e^{h_{\min}\sum\eps^3|\phi|^2}e^{\sum\eps^3(h-h_{\min})|\phi|^2}$ moves the wall into the bare coefficient, $\gamma'\defeq a_m-2h_{\min}$, and Jensen's inequality bounds the remaining factor below by $\exp[\sum\eps^3(h-h_{\min})\langle|\phi|^2\rangle_{\gamma'}]\ge1$.
\emph{Step 2 (the higher King point).} $u\defeq\frac12(a_M-\gamma')=\frac12(a_M-a_m+2h_{\min})$ lies in $[0,C]$ for $\eps\le\eps_0$, by $h_{\min}\ge-H$, $\Delta_0(M^2)>2H$ and the finite upper limit of $a_M-a_m$; since $-\partial_\gamma\ln\hat Z=\frac12\langle\sum\eps^3|\phi|^2\rangle_\gamma$ is nonincreasing in $\gamma$, $\ln\hat Z(\gamma')\ge\ln\hat Z(a_M)+u|T|(2c_M+w_k(M))$, where input (i) bounds $\ln\hat Z(a_M)=\ln\Zun^{M^2}(0,0)+P_M$ below by $-C|T|$ plus the $c$-numbers.
\emph{Step 3 ($c$-numbers).} By the integrated counterterm identity the four $\eps^{-1}$-divergent $c$-numbers collapse exactly to $2u(c_M-c_m)|T|=O(|T|)$; the remainder $|T|\int_{m^2}^{M^2}(c_m-c_s)Y\,ds$ is bounded below by $-C|T|$ from $0\le c_m-c_s\le C$ and $|Y|\le C$ (its sign is not needed), and the perturbative remainder is $\ge-C(1+k)|T|$ by (iii).
\emph{Step 4.} $w_k(M)\ge-C(1+k)$ is the mirror bound from (i)--(iii) at $M^2$, $M_+^2$: the same floor shift with the roles of the two points exchanged and the upper bound (ii) in place of the lower one.
\emph{Step 5.} Assemble; the source $tG$ is removed as in \Cref{rot:prop:H5}. No stability bound at $m^2$ itself, no bound on $w_k(m^2)$ and no two-sided Wick bound is used. The statement is in the free-boundary convention; the $c$-number mismatch with the torus convention is $\Delta c_T=O(e^{-c\ell(T)})$ (\Cref{rot:prop:CT}(b)).
\end{proof}
